% Bibliography of paper.tex and partA.tex.
\begin{thebibliography}{99}
\bibitem{lippert2002} R.~Lippert, R.~Schwartz, G.~Lancia, S.~Istrail. Algorithmic strategies for the single nucleotide polymorphism haplotype assembly problem. \emph{Briefings in Bioinformatics} 3(1):23--31, 2002. \url{https://doi.org/10.1093/bib/3.1.23}
\bibitem{greenberg2004} H.~J.~Greenberg, W.~E.~Hart, G.~Lancia. Opportunities for combinatorial optimization in computational biology. \emph{INFORMS J. Computing} 16(3):211--231, 2004. \url{https://doi.org/10.1287/ijoc.1040.0073}
\bibitem{patterson2015} M.~Patterson et al. WhatsHap: weighted haplotype assembly for future-generation sequencing reads. \emph{J. Comput. Biol.} 22(6):498--509, 2015. \url{https://doi.org/10.1089/cmb.2014.0157}
\bibitem{martin2016} M.~Martin et al. WhatsHap: fast and accurate read-based phasing. bioRxiv 085050, 2016. \url{https://doi.org/10.1101/085050}
\bibitem{fischer2016} S.~O.~Fischer, T.~Marschall. Selecting reads for haplotype assembly. bioRxiv 046771, 2016. \url{https://doi.org/10.1101/046771}
\bibitem{ralphi2025} E.~Battistella, A.~Maheshwari, B.~Ekim, B.~Berger, V.~Popic. Graph-based deep reinforcement learning for haplotype assembly with Ralphi. \emph{Genome Res.} 35(12):2617--2625, 2025. \url{https://doi.org/10.1101/gr.280569.125}
\bibitem{whatshap2015}
M.~Patterson, T.~Marschall, N.~Pisanti, L.~van Iersel, L.~Stougie,
G.~W.~Klau and A.~Sch\"onhuth.
\newblock WhatsHap: weighted haplotype assembly for future-generation
sequencing reads.
\newblock \emph{Journal of Computational Biology}, 22(6), 2015.
\newblock \doi{10.1089/cmb.2014.0157}.
\bibitem{whatshap2016}
M.~Martin et al.
\newblock WhatsHap: fast and accurate read-based phasing.
\newblock bioRxiv 085050, 2016.
\bibitem{hiphase}
J.~M.~Holt et al.
\newblock HiPhase: jointly phasing small, structural, and tandem repeat
variants from HiFi sequencing.
\newblock \emph{Bioinformatics}, 40(2):btae042, 2024.
\newblock \doi{10.1093/bioinformatics/btae042}.
\bibitem{longphase}
J.-H.~Lin et al.
\newblock LongPhase: an ultra-fast chromosome-scale phasing algorithm for
small and large variants.
\newblock \emph{Bioinformatics}, 38(7):1816--1822, 2022.
\newblock \doi{10.1093/bioinformatics/btac058}.
\bibitem{arkin1987} E.~M.~Arkin, E.~B.~Silverberg. Scheduling jobs with fixed start and end times. \emph{Discrete Appl. Math.} 18:1--8, 1987. \url{https://doi.org/10.1016/0166-218X(87)90037-0}
\bibitem{carlisle1995} M.~C.~Carlisle, E.~L.~Lloyd. On the k-coloring of intervals. \emph{Discrete Appl. Math.} 59:225--235, 1995. \url{https://doi.org/10.1016/0166-218X(95)80003-M}
\bibitem{richards2015}
S.~Richards et al.
\newblock Standards and guidelines for the interpretation of sequence
variants: a joint consensus recommendation of the American College of Medical
Genetics and Genomics and the Association for Molecular Pathology.
\newblock \emph{Genetics in Medicine}, 17(5):405--424, 2015.
\newblock \doi{10.1038/gim.2015.30}.
\bibitem{delgobbo2025}
G.~F.~Del~Gobbo and K.~M.~Boycott.
\newblock The additional diagnostic yield of long-read sequencing in
undiagnosed rare diseases.
\newblock \emph{Genome Research}, 35(4):559--571, 2025.
\newblock \doi{10.1101/gr.279970.124}.
\bibitem{shapeit5}
R.~J.~Hofmeister, D.~M.~Ribeiro, S.~Rubinacci and O.~Delaneau.
\newblock Accurate rare variant phasing of whole-genome and whole-exome
sequencing data in the UK Biobank.
\newblock \emph{Nature Genetics}, 55(7):1243--1249, 2023.
\newblock \doi{10.1038/s41588-023-01415-w}.
\bibitem{lassen2024}
F.~H.~Lassen et al.
\newblock Exome-wide evidence of compound heterozygous effects across common
phenotypes in the UK Biobank.
\newblock \emph{Cell Genomics}, 4(7):100602, 2024.
\newblock \doi{10.1016/j.xgen.2024.100602}.
\bibitem{geifman2017}
Y.~Geifman and R.~El-Yaniv.
\newblock Selective classification for deep neural networks.
\newblock In \emph{Advances in Neural Information Processing Systems 30}, 2017.
\bibitem{ltt2025}
A.~N.~Angelopoulos, S.~Bates, E.~J.~Cand\`es, M.~I.~Jordan and L.~Lei.
\newblock Learn then Test: calibrating predictive algorithms to achieve risk
control.
\newblock \emph{Annals of Applied Statistics}, 19(2), 2025.
\newblock \doi{10.1214/24-AOAS1998}.
\bibitem{dwarshuis2024} N.~Dwarshuis et al. The GIAB genomic stratifications resource for human reference genomes. \emph{Nat. Commun.} 15:9029, 2024. \url{https://doi.org/10.1038/s41467-024-53260-y}
\bibitem{vovk2012}
V.~Vovk.
\newblock Conditional validity of inductive conformal predictors.
\newblock Asian Conference on Machine Learning 2012, per arXiv.
\newblock arXiv:1209.2673.
\bibitem{he2010} D.~He, A.~Choi, K.~Pipatsrisawat, A.~Darwiche, E.~Eskin. Optimal algorithms for haplotype assembly from whole-genome sequence data. \emph{Bioinformatics} 26(12):i183--i190, 2010. \url{https://doi.org/10.1093/bioinformatics/btq215}
\bibitem{pirola2016} Y.~Pirola, S.~Zaccaria, R.~Dondi, G.~W.~Klau, N.~Pisanti, P.~Bonizzoni. HapCol: accurate and memory-efficient haplotype assembly from long reads. \emph{Bioinformatics} 32(11):1610--1617, 2016. \url{https://doi.org/10.1093/bioinformatics/btv495}
\bibitem{beretta2018} S.~Beretta et al. HapCHAT: adaptive haplotype assembly for efficiently leveraging high coverage in long reads. \emph{BMC Bioinformatics} 19:252, 2018. \url{https://doi.org/10.1186/s12859-018-2253-8}
\bibitem{chen2013} Z.-Z.~Chen, F.~Deng, L.~Wang. Exact algorithms for haplotype assembly from whole-genome sequence data. \emph{Bioinformatics} 29(16):1938--1945, 2013. \url{https://doi.org/10.1093/bioinformatics/btt349}
\bibitem{bonizzoni2016} P.~Bonizzoni, R.~Dondi, G.~W.~Klau, Y.~Pirola, N.~Pisanti, S.~Zaccaria. On the minimum error correction problem for haplotype assembly in diploid and polyploid genomes. \emph{J. Comput. Biol.} 23(9):718--736, 2016. \url{https://doi.org/10.1089/cmb.2015.0220}
\bibitem{edge2017} P.~Edge, V.~Bafna, V.~Bansal. HapCUT2: robust and accurate haplotype assembly for diverse sequencing technologies. \emph{Genome Res.} 27(5):801--812, 2017. \url{https://doi.org/10.1101/gr.213462.116}
\bibitem{duitama2012} J.~Duitama et al. Fosmid-based whole genome haplotyping of a HapMap trio child. \emph{Nucleic Acids Res.} 40(5):2041--2053, 2012. \url{https://doi.org/10.1093/nar/gkr1042}
\bibitem{das2015} S.~Das, H.~Vikalo. SDhaP: haplotype assembly for diploids and polyploids via semi-definite programming. \emph{BMC Genomics} 16:260, 2015. \url{https://doi.org/10.1186/s12864-015-1408-5}
\bibitem{makinen2015} V.~M\"akinen, V.~Staneva, A.~I.~Tomescu, D.~Valenzuela, S.~Wilzbach. Interval scheduling maximizing minimum coverage. \emph{Discrete Appl. Math.} 225:130--135, 2017. \url{https://doi.org/10.1016/j.dam.2016.08.021}
\bibitem{garg2018} S.~Garg. Computational haplotyping: theory and practice. PhD thesis, Saarland University, 2018.
\bibitem{bareyehuda2006} R.~Bar-Yehuda, M.~M.~Halld\'orsson, J.~Naor, H.~Shachnai, I.~Shapira. Scheduling split intervals. \emph{SIAM J. Comput.} 36(1):1--15, 2006. \url{https://doi.org/10.1137/S0097539703437843}
\bibitem{holt2024} J.~M.~Holt et al. HiPhase: jointly phasing small, structural, and tandem repeat variants from HiFi sequencing. \emph{Bioinformatics} 40(2):btae042, 2024. \url{https://doi.org/10.1093/bioinformatics/btae042}
\bibitem{lin2022} J.-H.~Lin et al. LongPhase: an ultra-fast chromosome-scale phasing algorithm for small and large variants. \emph{Bioinformatics} 38(7):1816--1822, 2022. \url{https://doi.org/10.1093/bioinformatics/btac058}
\bibitem{kamath2015} G.~M.~Kamath, E.~\c{S}a\c{s}o\u{g}lu, D.~Tse. Optimal haplotype assembly from high-throughput mate-pair reads. In \emph{Proc. IEEE Int. Symp. Inf. Theory (ISIT)}, 914--918, 2015. \url{https://doi.org/10.1109/ISIT.2015.7282588}
\bibitem{chen2016} Y.~Chen, G.~Kamath, C.~Suh, D.~Tse. Community recovery in graphs with locality. In \emph{Proc. 33rd Int. Conf. Mach. Learn. (ICML)}, PMLR 48:689--698, 2016. \url{https://proceedings.mlr.press/v48/chena16.html}
\bibitem{mahajan2024} D.~Mahajan, C.~Jain, N.~Kashyap. On the coverage required for diploid genome assembly. arXiv:2405.05734, 2024.
\bibitem{teo2024} A.~Y.~Y.~Teo, J.~W.~Squair, G.~Courtine, M.~A.~Skinnider. Best practices for differential accessibility analysis in single-cell epigenomics. \emph{Nat. Commun.} 15:8805, 2024. \url{https://doi.org/10.1038/s41467-024-53089-5}
\bibitem{mallona2025} I.~Mallona, C.~Soneson, B.~Carrillo, et al. Building a continuous benchmarking ecosystem in bioinformatics. \emph{PLoS Comput. Biol.} 21(11):e1013658, 2025. \url{https://doi.org/10.1371/journal.pcbi.1013658}
\bibitem{hapcut2}
P.~Edge, V.~Bafna and V.~Bansal.
\newblock HapCUT2: robust and accurate haplotype assembly for diverse
sequencing technologies.
\newblock \emph{Genome Research}, 27(5):801--812, 2017.
\newblock \doi{10.1101/gr.213462.116}.
\bibitem{probhap}
V.~Kuleshov.
\newblock Probabilistic single-individual haplotyping.
\newblock \emph{Bioinformatics}, 30(17), 2014.
\bibitem{damaraju2026} N.~Damaraju, F.~G.~Frost, J.~Fu, D.~A.~D'Onofrio, J.~Goffena, S.~H.~Storz, Z.~B.~Anderson, T.~Prall, M.~P.~Zalusky, M.~C.~V.~Malicdan, D.~R.~Adams, D.~E.~Miller. Determinants of haplotype phasing accuracy in long-read human genome sequencing. \emph{bioRxiv} preprint, 2026. \url{https://doi.org/10.64898/2026.05.04.722832}
\bibitem{phaseme}
M.~Majidian and F.~J.~Sedlazeck.
\newblock PhaseME: automatic rapid assessment of phasing quality and phasing
improvement.
\newblock \emph{GigaScience}, 9(7):giaa078, 2020.
\newblock \doi{10.1093/gigascience/giaa078}.
\bibitem{hiphap}
J.~Heinz, M.~Marin, M.~Meyerson and H.~Li.
\newblock HipHap: haplotype assignment and confidence scoring for diploid
reference genomes.
\newblock bioRxiv, posted 29 September 2026.
\newblock \doi{10.64898/2026.09.24.754153}.
\bibitem{chen2025paraphase} X.~Chen et al. Genome-wide profiling of highly similar paralogous genes using HiFi sequencing. \emph{Nat. Commun.} 16:2340, 2025. \url{https://doi.org/10.1038/s41467-025-57505-2}
\bibitem{jamshidi2025} Jamshidi et al. Long-range PCR and Nanopore sequencing for localisation and phasing variants: an end-to-end clinical application workflow. \emph{BMC Med. Genomics} 18:186, 2025. \url{https://doi.org/10.1186/s12920-025-02251-z}
\bibitem{devaney2026} Devaney et al. Sensitivity of HiFi long-read genome sequencing for difficult-to-detect pathogenic variants when applied to real-world clinical laboratory samples. \emph{Am. J. Hum. Genet.} 113(5):1036--1048, 2026. \url{https://doi.org/10.1016/j.ajhg.2026.04.001}
\bibitem{crc2024}
A.~N.~Angelopoulos, S.~Bates, A.~Fisch, L.~Lei and T.~Schuster.
\newblock Conformal risk control.
\newblock In \emph{International Conference on Learning Representations}, 2024.
\newblock arXiv:2208.02814.
\bibitem{gibbs2023}
Gibbs, Cherian and E.~J.~Cand\`es.
\newblock Conformal prediction with conditional guarantees.
\newblock arXiv:2305.12616, 2023.
\bibitem{tibshirani2019}
R.~J.~Tibshirani, R.~F.~Barber, E.~J.~Cand\`es and A.~Ramdas.
\newblock Conformal prediction under covariate shift.
\newblock arXiv:1904.06019, 2019.
\bibitem{hansen2026} N.~F.~Hansen, N.~Dwarshuis, H.~J.~Ji, A.~Rhie, et al.\ (65 authors), S.~Koren, J.~M.~Zook, A.~M.~Phillippy. A complete diploid human genome benchmark for personalized genomics. \emph{Cell} 189(16):4857--4875.e31, 2026. \url{https://doi.org/10.1016/j.cell.2026.06.016}
\bibitem{kronenberg2025}
Z.~Kronenberg et al.
\newblock The Platinum Pedigree: a long-read benchmark for genetic variants.
\newblock \emph{Nature Methods}, 22:1669--1676, 2025.
\newblock \doi{10.1038/s41592-025-02750-y}.
\bibitem{browning2022}
Browning and Browning.
\newblock Genotype error biases trio-based estimates of haplotype phase
accuracy.
\newblock \emph{American Journal of Human Genetics}, 109:1016--1025, 2022.
\newblock \doi{10.1016/j.ajhg.2022.04.019}.
\bibitem{kong2008} A.~Kong et al. Detection of sharing by descent, long-range phasing and haplotype imputation. \emph{Nat. Genet.} 40(9):1068--1075, 2008. \url{https://doi.org/10.1038/ng.216}
\bibitem{oconnell2016} J.~O'Connell et al. Haplotype estimation for biobank-scale data sets. \emph{Nat. Genet.} 48(7):817--820, 2016. \url{https://doi.org/10.1038/ng.3583}
\bibitem{akbari2022} V.~Akbari et al. Genome-wide detection of imprinted differentially methylated regions using nanopore sequencing. \emph{eLife} 11:e77898, 2022. \url{https://doi.org/10.7554/eLife.77898}
\bibitem{majidian2020} S.~Majidian, M.~H.~Kahaei, D.~de Ridder. Minimum error correction-based haplotype assembly: considerations for long read data. \emph{PLOS ONE} 15(6):e0234470, 2020. \url{https://doi.org/10.1371/journal.pone.0234470}
\bibitem{zook2016} J.~M.~Zook et al. Extensive sequencing of seven human genomes to characterize benchmark reference materials. \emph{Sci. Data} 3:160025, 2016. \url{https://doi.org/10.1038/sdata.2016.25}
% Part III references; titles, volumes and pages checked against Crossref on 18 October 2026.
\bibitem{wagner2022} J.~Wagner et al. Benchmarking challenging small variants with linked and long reads. \emph{Cell Genomics} 2(5):100128, 2022. \url{https://doi.org/10.1016/j.xgen.2022.100128}
\bibitem{berend2005} D.~Berend, L.~Sapir. Monotonicity in Condorcet jury theorem. \emph{Soc. Choice Welf.} 24(1):83--92, 2005. \url{https://doi.org/10.1007/s00355-003-0293-z}
\end{thebibliography}
